\section{Proof-Obligation Comparison}
\label{app:comparison}

The comparison below summarizes the forms of proof obligations used in the two FDR proofs. It is a qualitative comparison of their representation, not a measurement of proof-development or solver runtime.

\begin{center}
\begin{tabularx}{\linewidth}{@{}>{\raggedright\arraybackslash}p{0.25\linewidth}>{\raggedright\arraybackslash}X>{\raggedright\arraybackslash}X@{}}
\toprule
 & pCFG level & program level\\
\midrule
Invariant & location-indexed state predicate & one distributional invariant\\
Supermartingale & edge/choice obligations & one expected body-step inequality\\
Variant & edge-wise strict\slash probabilistic descent & one post-distribution inequality\\
Control-flow constants & scaling and location offsets & none\\
Reuse from partial correctness & requires projection/lifting & direct\\
\bottomrule
\end{tabularx}
\end{center}

For FDR, both columns are feasible and prove the same AST fact. This makes FDR a controlled comparison of proof style. The program-level proofs of FLDR and Han--Hoshi further demonstrate how progress is expressed using non-reset outcomes and interval boundaries, respectively; they are not additional comparisons against explicitly constructed pCFG certificates.
